\section{Rational caustics and analytic regularity}\label{ana:section}

The first task is to obtain analytic regularity from the continuous
foliation in \Cref{def:caustic-collar}.  The argument uses only the rational
leaves with rotation number \(1/q\).  Their critical actions define a system
of equations whose derivative on sufficiently high Fourier modes is close
to the identity.  The small loss of complex width in the construction of
each individual equation is compatible with an inverse on a single
exponentially weighted Fourier space.
The coordinate follows Lazutkin's grazing normalization, and the
action viewpoint is related to Marvizi--Melrose's asymptotic
interpolation~\cite{Lazutkin1973,MM1982}. Rational-caustic Fourier
constraints also appear in the perturbative work of Avila, De Simoi
and Kaloshin~\cite{ADK2016}. The arbitrary-input mesh, Fourier row,
and nonperturbative inversion required here are proved below.

\begin{theorem}[Analytic regularity]\label{thm:analytic}
Under the hypotheses of \Cref{thm:main}, the support function of the table
and its boundary are real analytic.
\end{theorem}

\subsection{The invariant graphs and their rational leaves}

Classical invariant-graph results for twist maps give context for this
phase description~\cite{Arnaud2009}; periodic billiard graphs are also
studied in \cite[Appendix~A]{FS2026}. The all-point rational periodicity
needed here follows directly from the filled collar in the next lemma.

Translate an interior point to the origin.  Write
\[
n_\theta=(\cos\theta,\sin\theta),\qquad
e_\theta=(-\sin\theta,\cos\theta).
\]
The support function \(h\) is smooth, and the point with outward normal
\(n_\psi\) is
\[
\gamma(\psi)=h(\psi)n_\psi+h'(\psi)e_\psi.
\]
The curvature assumption gives
\[
r(\psi)=h(\psi)+h''(\psi)>0.
\]
All angles below are real lifts of variables of period \(2\pi\).
An oriented line near a positively oriented supporting line is described
by its velocity \(e_\theta\) and its equation \(n_\theta\cdot x=p\).
If its forward reflection has outgoing angle \(\theta'\), put
\(\psi=(\theta+\theta')/2\) and \(d=(\theta'-\theta)/2\).  On the small
positive branch, the reflection law gives
\[
\begin{split}
\theta&=\psi-d,\qquad \theta'=\psi+d,\\
p&=h(\psi)\cos d-h'(\psi)\sin d,\\
p'&=h(\psi)\cos d+h'(\psi)\sin d.
\end{split}
\]
The billiard map has Bialy and Mironov's support-function generating
action~\cite{BM2017}:
\begin{equation}\label{ana:line-generating-function}
\mathcal S(\theta,\theta')
=2h\left(\frac{\theta+\theta'}2\right)
 \sin\left(\frac{\theta'-\theta}2\right),
\qquad
p=-\mathcal S_1,\quad p'=\mathcal S_2.
\end{equation}
In particular,
\begin{equation}\label{ana:twist}
\mathcal S_{12}(\theta,\theta')
=\frac{r(\psi)}2\sin d>0.
\end{equation}
This map preserves \(\dd p\wedge\dd\theta\), and
\(\partial\theta'/\partial p=-1/\mathcal S_{12}<0\).

\begin{lemma}\label{lem:rational-periodicity}
There is a one-sided phase collar filled by invariant continuous graphs
\(p=g_b(\theta)\), where \(b>0\) increases towards the interior and
\(g_b\) strictly decreases pointwise.  The induced increasing lifts
\(F_b\) have positive translation numbers \(\tau(b)\), tending to zero
as \(b\downarrow0\).

For every sufficiently small positive rational number \(l/q\), in lowest
terms, there is exactly one graph with
\(\tau(b)=2\pi l/q\), and its lift satisfies
\begin{equation}\label{ana:rational-identity}
F_b^q(\theta)=\theta+2\pi l
\qquad\text{for every }\theta.
\end{equation}
\end{lemma}

\begin{proof}
Restrict the foliation parameter to a compact interval close to zero.
Continuity of its parametrizing homeomorphism gives uniform convergence
of the leaf curves to the outer boundary.  Support functions are
continuous for this convergence.  Hence the convex bodies bounded by
the leaves contain the origin for sufficiently small parameters.
Two disjoint convex boundary curves surrounding the origin are nested:
otherwise their radial functions change order and the curves intersect.
Their compact disjointness makes the nesting strict, including strict
ordering of their support functions in every direction.  It follows,
using the limiting outer boundary, that \(g_b(\theta)\) is continuous
and strictly decreasing in \(b\), for each \(\theta\).
The intervals it fills, as \(\theta\) varies, form a phase collar.

Every positively oriented supporting line of a leaf in this collar
reflects to another positively oriented supporting line of the same
leaf.  Indeed the reflected line is tangent by the caustic property.
For a sufficiently small collar both the incoming and outgoing lines
are close to the corresponding positive outer supporting lines.
The other orientation of a leaf support has coordinate
\(-g_b(\theta+\pi)\), which stays negative because the bodies contain
a fixed neighborhood of the origin.  Thus this orientation cannot
occur.  The same argument applies to the inverse reflection.
The billiard map therefore restricts to a circle homeomorphism of each
graph.  It is close to the identity and consequently has degree one.
Its lift \(F_b\) is increasing, has strictly positive increments, and
depends continuously on \(b\).  By \Cref{ana:twist}, these lifts are
strictly ordered as \(b\) increases.

We recall the elementary rotation facts needed here.  For an increasing
lift \(F\) commuting with translation by \(2\pi\), the oscillation of
\(F^m(x)-x\) is at most \(2\pi\).  The maximum and minimum of this
displacement are respectively subadditive and superadditive in \(m\).
It follows that
\[
\tau(F)=\lim_{m\to\infty}\frac{F^m(x)-x}{m}
\]
exists, is independent of \(x\), and satisfies
\[
\left|\tau(F)-\frac{F^m(x)-x}{m}\right|\le\frac{2\pi}{m}.
\]
This bound gives continuity of \(\tau\) under uniform convergence of
lifts.  Ordering lifts orders their iterates and hence their translation
numbers.  A lift with translation number zero has a fixed point: if its
displacement never vanishes, its sign is constant and its absolute value
has a positive minimum.  Applied to \(F^q-2\pi l\), this says that
\(\tau(F)=2\pi l/q\) implies the existence of a \(q\)-periodic point.
Conversely, such a point forces that translation number.

These facts show that \(\tau(b)\) is continuous and nondecreasing.
Its positivity follows from the positive minimum of \(F_b(\theta)-\theta\)
on each fixed graph; its limit at zero follows from uniform convergence
of the map to the identity.

There cannot be a nontrivial interval of parameters with the same
rational translation number \(2\pi l/q\).  Otherwise the closed annulus
between two distinct graphs in that interval is invariant and has
positive finite area.  Almost every point in it is recurrent for the
\(q\)-th iterate, by finite-measure recurrence.  On a graph of the
specified rotation number, \(F_b^q-2\pi l\) has fixed points.  Every
point in a component of their complement moves strictly monotonically
towards an endpoint of that component and is not recurrent.
Thus the \(q\)-th iterate is the identity almost everywhere in the
annulus, and hence everywhere by continuity.  This contradicts strict
ordering of the \(q\)-th iterates on its two boundary graphs.

Continuity now supplies a unique graph at each sufficiently small
positive rational rotation, with smaller and larger rotations
arbitrarily close on its two sides.  A lift with rotation greater than
\(2\pi l/q\) has \(F^q(\theta)-\theta-2\pi l>0\) everywhere: a zero
would force that rational rotation, and a strictly negative displacement
would force a smaller one.  The corresponding inequality is reversed
below that rotation.  Passing to the graph in question proves
\Cref{ana:rational-identity}.
\end{proof}

\subsection{The normalized action}

Choose \(k>0\) so that the coordinate defined by
\[
\frac{\dd x}{\dd\theta}=k r(\theta)^{1/3}
\]
has period \(2\pi\).  Write \(\theta=\phi(x)\), and set
\begin{equation}\label{ana:normalization}
\phi'(x)=e^{u(x)},\qquad
r(\phi(x))=c e^{-3u(x)},\qquad c=k^{-3}.
\end{equation}
For the table, \(u\) is real, smooth, and periodic.  Subtract
\(\int_\theta^{\theta'}h(\eta)\,\dd\eta\) from \(\mathcal S\), change
coordinates, and divide by \(c\).  The subtracted term telescopes in
every fixed-endpoint action, so these operations preserve stationarity.
Integration by parts using \(h+h''=r\) gives the resulting action:
\begin{equation}\label{ana:normalized-action}
\begin{split}
D(X,Y)={}&-\int_X^\mu
 \bigl(1-\cos(\phi(z)-\phi(X))\bigr)e^{-2u(z)}\,\dd z\\
&-\int_\mu^Y
 \bigl(1-\cos(\phi(Y)-\phi(z))\bigr)e^{-2u(z)}\,\dd z,
\qquad
\phi(\mu)=\frac{\phi(X)+\phi(Y)}2 .
\end{split}
\end{equation}
For clarity, before the coordinate change the corresponding two
integrands are \(r(\eta)(1-\cos(\eta-\theta))\) and
\(r(\eta)(1-\cos(\theta'-\eta))\), respectively.  Substitution of
\(r(\phi(z))\phi'(z)/c=e^{-2u(z)}\) gives
\Cref{ana:normalized-action}.

The following action formula and mesh estimates apply to arbitrary
input; the caustic hypothesis is used separately to make the resulting
constraints vanish. We will use this formula for arbitrary real smooth
periodic \(u\).
In that case let \(\phi\) be any primitive of \(e^u\).  Its total
increment need not be \(2\pi\): all occurrences of \(\phi\) are
differences, so \(D\) remains diagonally \(2\pi\)-periodic.  For small
\(Y-X\), the midpoint is uniquely defined.  The formula also gives
\begin{equation}\label{ana:action-symmetries}
D(X+2\pi,Y+2\pi)=D(X,Y),\qquad D(Y,X)=-D(X,Y).
\end{equation}
If \(u\) is holomorphic in a strip, these definitions are holomorphic
near its diagonal.  All small complex integrals may be taken along
straight segments.

Differentiating the two integrals in \Cref{ana:normalized-action}, the
moving-midpoint terms cancel.  Thus
\begin{equation}\label{ana:D1-integral}
D_1(X,X+T)=
e^{u(X)}\int_X^\mu
\sin(\phi(z)-\phi(X))e^{-2u(z)}\,\dd z.
\end{equation}
The following expansion will also fix the normalization used in the
next Section:
\begin{equation}\label{ana:action-expansion}
\begin{split}
D_1(X,X+T)&=\frac{T^2}{8}+a(X)T^4+O(T^5),\\
a(X)&=\frac{7u''(X)-8u'(X)^2-e^{2u(X)}}{384},\\
D(X,X+T)&=-\frac{T^3}{24}-\frac{a(X)}5T^5+O(T^6).
\end{split}
\end{equation}
Here and below the remainders have the differentiated bounds specified
in the lemmas where they are used.  One direct calculation is as
follows.  Put \(\alpha=u'(X)\), \(\beta=u''(X)\), \(A=e^{u(X)}\), and
\(z=(\mu-X)/T\).  The midpoint equation gives
\[
z=\frac12+\frac{\alpha T}{8}+\frac{\beta T^2}{16}+O(T^3).
\]
After rescaling \Cref{ana:D1-integral}, its integrand is
\[
v-\frac32\alpha T v^2
+\frac{7\alpha^2-5\beta-A^2}{6}T^2v^3+O(T^3).
\]
Integrating to \(z\) gives the first two lines of
\Cref{ana:action-expansion}; in particular its linear term in \(T\)
vanishes.  The leading term of \(D\) follows directly from its integral.
Antisymmetry forces its fourth-order coefficient to vanish.
Finally, if \(d_5(X)\) is its fifth-order coefficient, differentiation
with respect to the first endpoint at fixed second endpoint gives
\(-5d_5=a\), proving the last line.

\subsection{Uniform stationary meshes}

For \(\sigma>0\) let
\[
\mathcal A_\sigma=\{z\in\C:|\Im z|<\sigma\},\qquad
\|f\|_{a,\sigma}=\max_{0\le j\le a}
                  \sup_{\mathcal A_\sigma}|f^{(j)}|.
\]
Functions in a positive strip are holomorphic and periodic; they are
real symmetric if \(f(\bar z)=\overline{f(z)}\).
For \(\sigma=0\), the same norm means the \(C^a\) norm on the real
circle.  Constants in the lemmas below do not depend on \(\sigma\).

\begin{lemma}\label{ana:mesh-lemma}
Fix an integer \(m\ge2\), a bound \(M\), and put
\(p=m+2\), \(r_0=m+15\).
There are \(R,K,N\), depending only on \(M,m\), with the following
properties.  Suppose \(0\le\sigma\le1\), \(u\) is real symmetric, and
\(\|u\|_{r_0,\sigma}\le M\).  For every integer \(q\ge N\), let
\(\delta=2\pi/q\).  There is a stationary fixed-endpoint mesh
\begin{equation}\label{ana:mesh-definition}
X_j(x)=x+j\delta+w_j(x),\qquad 0\le j\le q,\qquad
w_0=w_q=0,
\end{equation}
on \(\mathcal A_{\sigma'}\), where
\(\sigma'=(1-K/q^2)\sigma\), satisfying
\begin{equation}\label{ana:mesh-norm}
\|w\|_*:=
q^2\max_j\|w_j\|_{p,\sigma'}
+q^3\max_j\|w_{j+1}-w_j\|_{p,\sigma'}\le R.
\end{equation}
It depends differentiably on smooth real variations of \(u\), and
\[
\|\dot w\|_*\le C\|\dot u\|_{r_0,\sigma}.
\]
The analogous difference estimate holds for two inputs in this ball.
On the real axis, the mesh is also unique pointwise in the ball with
the same scaled position and step bounds.  This real uniqueness ball
may be enlarged by any fixed amount if \(N\) is increased.
\end{lemma}

\begin{proof}
We give the composition estimates, including their width dependence.
For a trial mesh in the ball \Cref{ana:mesh-norm}, set
\[
X=X_j,\qquad s=X_j-X_{j-1},\qquad t=X_{j+1}-X_j.
\]
Then \(s,t=\delta+O(q^{-3})\) in \(C^p\).
For \(T=t\) or \(T=-s\), and \(0\le\lambda\le1\), write
\(\mu=X+\lambda T z_\lambda\).  The midpoint equation is
\begin{equation}\label{ana:midpoint-equation}
\int_0^{z_\lambda}
 e^{u(X+\lambda Tv)-u(X)}\,\dd v
=\frac12\int_0^1e^{u(X+\lambda Tv)-u(X)}\,\dd v .
\end{equation}
Moving the integrals of the exponential minus one to the other side
expresses \(z_\lambda-1/2\) as a contraction on a \(C^p\) ball of
radius \(C/q\).  Indeed the exponential minus one is \(O(q^{-1})\);
its derivative with respect to \(z_\lambda\), after writing the
left integral as an integral over \(v=\alpha z_\lambda\), has the
same bound.  Products and compositions in this argument are estimated
by the ordinary finite-order chain rule.

All its evaluation points lie in the original strip.  Their maps from
\(x\) are real symmetric and have derivative \(1+O(q^{-2})\):
\[
X'=1+O(q^{-2}),\qquad T'=O(q^{-3}),\qquad
z_\lambda'=O(q^{-1}).
\]
For a real-symmetric map \(x+\rho+e(x)\), with \(\rho\) real and
\(\|e'\|\le Cq^{-2}\), integration along the vertical segment from
\(\Re x\) gives
\[
|\Im(x+\rho+e(x))|\le(1+Cq^{-2})|\Im x|.
\]
Taking \(K\) sufficiently large therefore puts every such point,
and every segment used in the integrals, inside \(\mathcal A_\sigma\).
This is a relative loss \(K\sigma/q^2\), with no reciprocal power of
\(\sigma\).

Implicit differentiation of \Cref{ana:midpoint-equation} gives, for
\(1\le i\le3\),
\begin{equation}\label{ana:midpoint-jets}
\|\partial_\lambda^i z_\lambda\|_{p,\sigma'}\le Cq^{-i}.
\end{equation}
To see the scaling, the coefficient of the derivative being solved
for is \(e^{u(X+\lambda Tz_\lambda)-u(X)}=1+O(q^{-1})\).
Each explicit \(\lambda\)-derivative contributes a factor \(T\);
the lower implicit derivatives satisfy the same induction.
For a mesh variation,
\[
\|\dot X\|_{p,\sigma'}\le q^{-2}\|\dot w\|_*,
\qquad
\|\dot T\|_{p,\sigma'}\le q^{-3}\|\dot w\|_*.
\]
Differentiating the same equations therefore gives, for \(1\le i\le3\),
\begin{equation}\label{ana:midpoint-variation}
\|\partial_\lambda^i\dot z_\lambda\|_{p,\sigma'}
\le Cq^{-i}
 \bigl(q^{-2}\|\dot w\|_*+\|\dot u\|_{r_0,\sigma}\bigr).
\end{equation}
For \(i=0\) the right hand side has an additional factor \(q^{-1}\).
For example, a difference \(u(B)-u(X)\) is written as the integral
of \(u'\) from \(X\) to \(B\).  Its variation in a changed basepoint
uses a difference of \(u'\), hence retains the factor \(B-X=O(q^{-1})\).
This proves the extra factor also for variations, rather than merely
for the values.

The rescaled expression for \Cref{ana:D1-integral} is
\[
\begin{split}
F(\lambda)
&=(\lambda T)^{-2}D_1(X,X+\lambda T)\\
&=e^{u(X)}\int_0^{z_\lambda}
\frac{\sin\left(\lambda T\int_0^v
                 e^{u(X+\lambda Tv')}\,\dd v'\right)}{\lambda T}
 e^{-2u(X+\lambda Tv)}\,\dd v .
\end{split}
\]
The quotient at zero is removable: write it as
\(A_\lambda(v)\operatorname{sinc}(\lambda T A_\lambda(v))\), where
\(A_\lambda(v)=\int_0^v e^{u(X+\lambda Tv')}\,\dd v'\) and
\(\operatorname{sinc}z=\sin z/z\).
Thus \Cref{ana:midpoint-jets,ana:midpoint-variation} apply to this
expression and its first variations.  The calculation following
\Cref{ana:action-expansion} gives
\[
F(0)=\frac18,\qquad F'(0)=0,\qquad F''(0)=2a(X)T^2.
\]
More explicitly, its third derivative has \(C^p\) norm at most
\(Cq^{-3}\), and the third derivative of its first variation has
norm at most
\[
Cq^{-3}\bigl(q^{-2}\|\dot w\|_*+\|\dot u\|_{r_0,\sigma}\bigr).
\]
Taylor's integral remainder gives
\begin{equation}\label{ana:D1-uniform-expansion}
D_1(X,X+T)
=T^2\left(\frac18+a(X)T^2+T^3b(X,T)\right),
\end{equation}
where the evaluated functions \(a(X)\), \(b(X,T)\) have bounded
\(C^p\) norms.  Their first variations are bounded by
\[
C\bigl(q^{-2}\|\dot w\|_*+\|\dot u\|_{r_0,\sigma}\bigr).
\]
In dividing the Taylor remainder by \(T^3\), no small denominator
is lost: \(qT\) is \(C^p\)-close to \(2\pi\) or \(-2\pi\), while
\(\dot T/T=O(q^{-2}\|\dot w\|_*)\).

We record a finite derivative budget for these statements.  Third
\(\lambda\)-derivatives of the midpoint and integrand use derivatives
of \(u\) through order \(p+3\); their first variations use order \(p+4\).
The estimates for differences follow from
\[
f\circ A-f\circ B
=\int_0^1 f'(B+\alpha(A-B))(A-B)\,\dd\alpha
\]
and the \(C^p\) chain rule, with the straight segment inside the
input strip.  Allowing two further derivatives covers all the
coefficient and composition bounds just used.  Thus \(p+6\)
derivatives already suffice here, fewer than \(r_0=p+13\).
There has been no Cauchy estimate on \(u\), so these constants are
uniform also as \(\sigma\downarrow0\).

Antisymmetry of \(D\) makes the Euler equation at \(X_j\)
\[
D_1(X,X+t)-D_1(X,X-s)=0.
\]
Using \Cref{ana:D1-uniform-expansion} and factoring \(t^2-s^2\)
rewrites it as
\begin{equation}\label{ana:mesh-fixed-point}
w_{j+1}-2w_j+w_{j-1}
=-\frac{t^5b(X,t)+s^5b(X,-s)}
 {(t+s)\bigl(1/8+(t^2+s^2)a(X)\bigr)}.
\end{equation}
The right hand side is \(O(q^{-4})\) in \(C^p\), and its Lipschitz
constant in the starred mesh norm is \(O(q^{-6})\).
For example its numerator is \(O(q^{-5})\), its denominator has
size comparable to \(q^{-1}\), and a unit starred variation changes
each coefficient and each relative step by \(O(q^{-2})\).
These statements hold for complex values as well, since the
denominator is uniformly separated from zero.

The Dirichlet inverse of the second difference has norm at most
\(Cq^4\) from uniform \(C^p\) forcing into the starred norm.
Indeed, if the forcing is \(f_j\), set \(d_j=w_j-w_{j-1}\).
Then
\[
d_{j+1}=d_j+f_j,\qquad
d_1=-\frac1q\sum_{i=1}^{q-1}(q-i)f_i.
\]
It follows that \(\max\|d_j\|_{C^p}\le Cq\max\|f_j\|_{C^p}\)
and \(\max\|w_j\|_{C^p}\le Cq^2\max\|f_j\|_{C^p}\).
Equation \eqref{ana:mesh-fixed-point} is consequently a contraction
with constant \(Cq^{-2}\).  Choose \(R\) larger than its uniform
image bound, then choose the width constant and \(N\).
The image bound can be taken independently of this fixed \(R\):
once \(q\) is large depending on \(R\), all scaled steps and
derivatives in the preceding estimates lie in fixed bounded sets.

Differentiating the contraction gives the asserted input derivative
bound; comparing its fixed points gives the corresponding Lipschitz
bound.  The same argument with order-zero real norms proves pointwise
uniqueness and works in any larger fixed real ball after increasing
\(N\).  It follows in particular that constructions made with different
fixed smoothness bounds or widths agree on their real domains for all
sufficiently large \(q\).
\end{proof}

\subsection{The high-frequency equations}

For the mesh in \Cref{ana:mesh-lemma}, define its critical action and
two associated Fourier coefficients by
\begin{equation}\label{ana:constraint-definition}
W_q(x;u)=\sum_{j=0}^{q-1}D(X_j(x),X_{j+1}(x)),\qquad
C_n(u)=-\frac{q^2}{6\pi}[W_q(\,\cdot\,;u)]_n,\quad n=\pm q,
\end{equation}
where
\[
[f]_n=\frac1{2\pi}\int_0^{2\pi}e^{-\ii nx}f(x)\,\dd x.
\]
Real pointwise uniqueness makes these definitions compatible when
the construction is performed on different strips or with different
fixed smoothness bounds.

\begin{proposition}\label{prop:high-frequency-estimate}
For each integer \(m\ge2\) and \(M<\infty\), take \(r_0=m+15\).
There are constants \(K,N\), independent of \(0\le\sigma\le1\),
such that, whenever \(\|u\|_{r_0,\sigma}\le M\) and \(|n|\ge N\),
the derivative of \Cref{ana:constraint-definition} has the form
\begin{equation}\label{ana:constraint-derivative}
(\dd C_n)_u(v)=[b_n(u)v]_n.
\end{equation}
Here
\begin{equation}\label{ana:coefficient-estimate}
\|b_n(u)-1\|_{m,(1-K/n^2)\sigma}\le\frac{K}{|n|},
\end{equation}
and, for two such inputs on the same strip,
\begin{equation}\label{ana:coefficient-lipschitz}
\|b_n(u)-b_n(\widetilde u)\|_{m,(1-K/n^2)\sigma}
\le\frac{K}{|n|}\|u-\widetilde u\|_{r_0,\sigma}.
\end{equation}
The derivative identity holds for smooth real variations, and for
real-symmetric holomorphic variations with the indicated finite norms.
\end{proposition}

\begin{proof}
Put \(q=|n|\) and \(p=m+2\).  In varying the critical action, all interior
mesh-variation terms vanish by stationarity.  The endpoints \(x\)
and \(x+2\pi\) are fixed.  We may therefore keep the vertices fixed
and differentiate \(D\) directly with respect to \(u\).

On the left half of a cell \((X,X+t)\), differentiation of
\Cref{ana:normalized-action} and reversal of the two integrations give
the kernel
\begin{equation}\label{ana:variation-kernel}
\begin{split}
\mathcal K_-(z)={}&
2\bigl(1-\cos(\phi(z)-\phi(X))\bigr)e^{-2u(z)}\\
&-e^{u(z)}
\int_z^\mu\sin(\phi(\zeta)-\phi(X))e^{-2u(\zeta)}\,\dd\zeta ,
\end{split}
\end{equation}
integrated against \(v(z)\,\dd z\).
On the right half, the corresponding formula is
\[
\begin{split}
\mathcal K_+(z)={}&
2\bigl(1-\cos(\phi(X+t)-\phi(z))\bigr)e^{-2u(z)}\\
&-e^{u(z)}
\int_\mu^z\sin(\phi(X+t)-\phi(\zeta))
                  e^{-2u(\zeta)}\,\dd\zeta .
\end{split}
\]
The moving-midpoint terms cancel because the two original integrands
agree at \(\mu\).

Parametrize the two halves, always from left to right, by
\[
z_-(x,\alpha)=X+\alpha(\mu-X),\qquad
z_+(x,\alpha)=\mu+\alpha(X+t-\mu),\qquad 0\le\alpha\le1,
\]
and set
\(\tau_-=\alpha/2\), \(\tau_+=(1+\alpha)/2\).
Uniformly in \(j,\alpha\), the mesh and midpoint estimates give
\begin{equation}\label{ana:half-cell-estimates}
\begin{split}
z_\pm(x,\alpha)
 &=x+(j+\tau_\pm)\delta+O_{C^p}(q^{-2}),\\
(z_\pm)_\alpha&=\frac{\pi}{q}+O_{C^p}(q^{-2}),\\
\mathcal K_\pm(z_\pm(x,\alpha))
 &=q^{-2}k_0(2\pi\tau_\pm)+O_{C^p}(q^{-3}),
\end{split}
\end{equation}
where
\begin{equation}\label{ana:leading-kernel}
k_0(\xi)=\frac32\min(\xi,2\pi-\xi)^2-\frac{\pi^2}{2},
\qquad 0\le\xi\le2\pi.
\end{equation}
Each remainder in \Cref{ana:half-cell-estimates} has the corresponding
Lipschitz bound, with the same power of \(q\), in
\(\|u-\widetilde u\|_{r_0,\sigma}\).

Here is a verification of both the leading terms and the remainder
orders.  On a cell, replace each occurrence of \(u\) by its value at
\(X\), using the integrated identity for differences of \(u\).
Then replace \(\sin a\) by \(a\) and \(1-\cos a\) by \(a^2/2\);
every phase difference is \(O(q^{-1})\).  On the left this gives
\[
\mathcal K_-(z)
=(z-X)^2-\int_z^\mu(\zeta-X)\,\dd\zeta+O_{C^p}(q^{-3})
=\frac32(z-X)^2-\frac12(\mu-X)^2+O_{C^p}(q^{-3}).
\]
The constant exponential factors cancel.  The reflected expression
gives the right half.  Since
\[
z-X=\delta\tau_-+O_{C^p}(q^{-2}),\qquad
\mu-X=\delta/2+O_{C^p}(q^{-2}),
\]
and similarly from the right endpoint, these formulas give
\Cref{ana:leading-kernel}.  An integrated difference gains one
factor \(q^{-1}\) also after \(p\) derivatives.  Applying the
composition-difference identity used in \Cref{ana:mesh-lemma} proves
the same gain for input variations.  For completeness, the mesh
coefficient estimates used at most \(p+6\) input derivatives; one
further derivative for an integrated spatial difference and one for
its moving-point variation bound all the kernel estimates by
\(C^{p+8}=C^{m+10}\) input norms.  This is still below \(r_0\).

For each \(j,\alpha\) and sign, the real map \(x\mapsto z_\pm(x,\alpha)\)
is an increasing lift close to a real translation.  Write its inverse
as \(x=x_\pm(z,\alpha)\).
On a strip \(\mathcal A_{\sigma''}\), with
\(\sigma''=(1-K/q^2)\sigma\) and a larger \(K\), it satisfies
\begin{equation}\label{ana:inverse-half-cell}
x_\pm(z,\alpha)
=z-(j+\tau_\pm)\delta+O_{C^{m+1}}(q^{-2}).
\end{equation}
This estimate and its Lipschitz version use no additional input
derivatives.  In fact the forward error is \(O(q^{-2})\) in
\(C^{m+2}\).  Its inverse equation is a contraction in \(C^{m+1}\),
since a composition difference there uses just one additional
derivative of the forward error.  Real symmetry gives exactly the
relative imaginary displacement used in the mesh construction, so a
larger relative strip loss suffices for this inverse as well.

Change the Fourier integration variable from \(x\) to \(z\).
The integration interval may start at a different point, but its
length is \(2\pi\), and all integrands are periodic.  Formula
\eqref{ana:constraint-derivative} follows with
\begin{equation}\label{ana:bn-formula}
b_n(z)=
-\frac{q^2}{6\pi}
\sum_{j=0}^{q-1}\sum_{\pm}
\int_0^1
\mathcal K_\pm(z)\,
(z_\pm)_\alpha(x_\pm(z,\alpha),\alpha)\,
(x_\pm)_z(z,\alpha)\,
e^{-\ii n(x_\pm(z,\alpha)-z)}\,\dd\alpha .
\end{equation}
In this formula the kernel is that of the cell belonging to the
indicated \(j,\alpha,\pm\), evaluated after the inverse substitution.

The phase in \Cref{ana:bn-formula} has leading value
\[
e^{\ii n(j+\tau_\pm)\delta}
=e^{2\pi\ii\sgn(n)\tau_\pm}.
\]
Its error, and the error in each of its first \(m\) derivatives,
is \(O(q^{-1})\), since \(|n|=q\) multiplies the
\(O_{C^{m+1}}(q^{-2})\) inverse displacement.
The same statement holds for input differences.
Substituting the leading terms in \Cref{ana:bn-formula} thus gives
\[
-\frac13\int_0^1 k_0(2\pi\tau)e^{2\pi\ii\sgn(n)\tau}\,\dd\tau .
\]
The integral equals \(-3\).  Indeed symmetry reduces it to
\[
\frac1\pi\int_0^\pi
\left(\frac32\xi^2-\frac{\pi^2}{2}\right)\cos\xi\,\dd\xi=-3.
\]
The leading coefficient is therefore \(1\).

Each cell contributes \(O(q^{-1})\) to the leading coefficient and
\(O(q^{-2})\) to its error: the kernel, prefactor, and cell Jacobian
have respective sizes \(q^{-2},q^2,q^{-1}\), and their errors have
one further factor \(q^{-1}\). These are absolute estimates, including
at zeros of the leading kernel \(k_0\). Summation over \(q\) cells proves
\Cref{ana:coefficient-estimate}.  Since the leading term is independent
of \(u\), the same computation with differences proves
\Cref{ana:coefficient-lipschitz}.  The finite derivative budgets above
and the relative domain losses prove the stated uniformity in
\(\sigma\), including \(\sigma=0\).
\end{proof}

\begin{lemma}\label{ana:real-periodic-mesh}
For the function \(u\) associated with the actual table by
\Cref{ana:normalization},
\begin{equation}\label{ana:constraint-vanishing}
C_n(u)=0
\qquad\text{for all sufficiently large }|n|.
\end{equation}
\end{lemma}

\begin{proof}
The rational graphs of rotation \(2\pi/q\) tend to the boundary:
otherwise their parameters have a positive subsequential lower bound,
whereas every fixed positive leaf has positive rotation.
By \Cref{lem:rational-periodicity}, the successive \(x\)-coordinates
on each such graph form a positive stationary mesh of \(q\) steps
whose sum is \(2\pi\).  All its steps tend uniformly to zero.

Let \(s,t\) be consecutive steps.  The real Taylor expansion
\eqref{ana:D1-uniform-expansion} first gives \(s\asymp t\), and then,
after factoring \(t^2-s^2\), gives \(t-s=O(s^4)\).
Consequently
\[
|\log t-\log s|\le Cs^3.
\]
Summing along any part of a circuit costs at most
\[
C\sum_j s_j^3\le C(\max_j s_j)^2\sum_j s_j
=2\pi C(\max_j s_j)^2.
\]
All steps are therefore uniformly comparable.  Their sum implies
each is \(O(q^{-1})\), so adjacent differences are \(O(q^{-4})\).
It follows that
\[
s_j=\frac{2\pi}{q}+O(q^{-3}),\qquad
X_j-x-\frac{2\pi j}{q}=O(q^{-2}).
\]
These are exactly the real position and step bounds in a sufficiently
large fixed uniqueness ball of \Cref{ana:mesh-lemma}.
Thus this geometric mesh is the mesh defining \(W_q\), pointwise
for every initial \(x\).  In particular it has the smooth dependence
on \(x\) supplied by that construction.

The geometric orbit satisfies the Euler equation also at the glued
endpoints.  In differentiating \(W_q(x)\), its interior terms vanish
by stationarity and its two endpoint derivatives sum to zero.
Therefore \(W_q'(x)=0\), proving \Cref{ana:constraint-vanishing}.
\end{proof}

\subsection{A Wiener-space inverse without loss of width}

For a periodic function or Fourier sequence, put
\[
|f|_\sigma=\sum_{k\in\Z}|[f]_k|e^{\sigma|k|},
\qquad
|f|_{\sigma,A}=\sum_{k\in\Z}(1+|k|)^A|[f]_k|e^{\sigma|k|}.
\]
Let \(\Pi_{\ge N}\) denote projection onto indices \(|k|\ge N\),
and restrict the constraint map to those output indices.
The following estimate explains why the mode-dependent width in
\Cref{ana:coefficient-estimate} causes no loss for this operator.

\begin{lemma}\label{ana:wiener-estimate}
Fix the bounds in \Cref{prop:high-frequency-estimate} for \(m=2\),
and increase \(N\) if necessary.  On the Wiener space of width
\(\sigma\), the coefficient formula for \(\dd C(u)\) defines a
bounded operator \(A(u)\) satisfying
\begin{equation}\label{ana:wiener-near-identity}
A(u)=\Pi_{\ge N}+R(u),\qquad \|R(u)\|\le\frac{C}{N}.
\end{equation}
For two bounded inputs on the same strip,
\begin{equation}\label{ana:wiener-lipschitz}
\|A(u)-A(\widetilde u)\|
\le C\|u-\widetilde u\|_{r_0,\sigma}.
\end{equation}
In particular, the restriction \(A_H(u)\) to high input modes has
inverse of norm at most \(2\) when \(N\) is sufficiently large.
All these operators preserve real data.
\end{lemma}

\begin{proof}
Contour displacement and two integrations by parts give
\[
|[b_n(u)-1]_j|
\le\frac{C}{|n|}(1+|j|)^{-2}
 e^{-(1-K/n^2)\sigma|j|}.
\]
At \(\sigma=0\) only the integrations by parts are used.
For a positive strip one first displaces to a smaller closed strip
and then lets its width tend to the stated width.
The same estimate for coefficient differences has the additional
factor \(\|u-\widetilde u\|_{r_0,\sigma}\).

Take \(N^2\ge 2K\).  For all \(k\in\Z\) and \(|n|\ge N\),
\begin{equation}\label{ana:exponential-balance}
|n|-|k|-(1-K/n^2)|n-k|\le\frac{2K}{|n|}.
\end{equation}
If \(|k|\le|n|\), the triangle inequality leaves at most
\((K/n^2)|n-k|\le2K/|n|\).  If \(|k|>|n|\), the entire left hand
side is nonpositive.
The weighted matrix entry for \(R(u)\) is consequently bounded by
\[
\frac{C}{|n|}(1+|n-k|)^{-2}e^{2K\sigma/|n|}.
\]
Its column sum is at most
\[
\frac{C'}N\sum_{j\in\Z}(1+|j|)^{-2}\le\frac{C''}N.
\]
This proves \Cref{ana:wiener-near-identity}.  The difference estimate
follows from the same calculation.  A Neumann series gives the
high-mode inverse.  Finally \(C_{-q}(u)=\overline{C_q(u)}\) for real
inputs, so the operators and their inverse preserve real Fourier
symmetry.
\end{proof}

The operator in this lemma is initially obtained from smooth
variations and is then extended by the displayed bounded matrix.
This distinction matters at a strip boundary, where an element of
the Wiener space need not have \(r_0\) bounded derivatives.
On every smaller strip it does have these derivatives, and the
matrix coincides with the actual derivative there, by approximation
with Fourier polynomials.  It also coincides across input strip
widths: real pointwise mesh uniqueness identifies the critical actions
and hence all their derivative Fourier coefficients.

\subsection{Analytic approximation and uniqueness}

\begin{proof}[Proof of \Cref{thm:analytic}]
Let \(u\) be the actual smooth function in \Cref{ana:normalization}.
For the remainder of the proof set \(m=2\) and \(r_0=m+15\).
Choose
\[
M>4e\bigl(1+|u|_{0,r_0}\bigr),
\]
and choose \(N\) large enough for
\Cref{ana:constraint-vanishing,ana:wiener-estimate}, with
\(\|R\|\le1/2\).
These choices are fixed before the approximation parameter below.

Let \(u_0\) be the Fourier truncation of \(u\) to \(|k|\le P\),
where \(P>N\), and put \(\sigma_0=1/P\).
For each fixed integer \(a\),
\begin{equation}\label{ana:truncation-bound}
|u_0|_{\sigma_0,a}\le e|u|_{0,a}.
\end{equation}
We first prove that its residual is arbitrarily small in powers of
\(P\):
\begin{equation}\label{ana:initial-residual}
|C(u_0)|_{\sigma_0}=O(P^{-A})
\quad\text{for every fixed }A.
\end{equation}
The constants and the lower threshold for \(P\) may depend on \(A\).

For \(N\le|n|\le P\), the exponential weight is at most \(e\).
Integrate the derivative on the real segment from \(u\) to \(u_0\),
use \(C(u)=0\), and apply \Cref{ana:wiener-estimate} at width zero.
The sum of these weighted coefficients is bounded by
\[
C|u_0-u|_0=O(P^{-A})
\]
for every \(A\), since \(u\) is smooth.
For the remaining tail, \(C(0)=0\): the zero input has a uniform mesh
and a constant critical action.  Thus
\[
C_n(u_0)=\int_0^1(\dd C_n)_{\lambda u_0}(u_0)\,\dd\lambda.
\]
Use \Cref{prop:high-frequency-estimate} here with one fixed
integer \(m>A+2\).  By \Cref{ana:truncation-bound}, its higher input
norms are uniformly bounded in \(P\).  Once \(P\) exceeds the
corresponding threshold, the estimate applies to every \(|n|>P\).
Real mesh uniqueness identifies these constraints with the ones
already chosen using \(m=2\).  The weighted convolution estimate
in the preceding lemma also holds with a factor \((1+|n|)^A\):
use
\[
(1+|n|)^A\le
(1+|n-k|)^A(1+|k|)^A
\]
and sum the resulting power \((1+|n-k|)^{A-m}\).
It follows that
\[
\sum_{|n|>P}(1+|n|)^A|C_n(u_0)|e^{\sigma_0|n|}
\le C_A|u_0|_{\sigma_0,A}\le C'_A.
\]
Dividing by \((1+P)^A\) proves the tail bound and hence
\Cref{ana:initial-residual}.  Only one larger fixed smoothness
order is used for each requested power; no threshold uniform over
all smoothness orders is needed.

We now solve the high-mode equations while keeping every mode
\(|k|<N\) fixed.  Set
\[
\sigma_j=\frac{1+2^{-j}}{2P},
\qquad
\sigma_\infty=\frac1{2P},
\qquad
\eps_j=|C(u_j)|_{\sigma_j}.
\]
At step \(j\), let
\begin{equation}\label{ana:newton-correction}
\eta_j=-A_H(u_j)^{-1}C(u_j),\qquad u_{j+1}=u_j+\eta_j.
\end{equation}
The correction has only high modes and satisfies
\[
|\eta_j|_{\sigma_j}\le2\eps_j.
\]
For \(a=\sigma_j-\sigma_{j+1}=2^{-j-2}/P\), the elementary inequality
\(\sup_{\xi\ge0}\xi^{r_0}e^{-a\xi}\le C_{r_0}a^{-r_0}\) gives
\begin{equation}\label{ana:correction-smoothing}
\|\eta_j\|_{r_0,\sigma_{j+1}}
\le C P^{r_0}2^{r_0j}\eps_j.
\end{equation}
Thus the correcting segment has the finite input norm required for
the derivative estimate on the smaller strip.

As long as the inputs on that segment stay within the fixed bound
\(M\), integrate the difference of derivatives on it.
The linear term cancels exactly by
\Cref{ana:newton-correction}.  Although its inverse was computed at
width \(\sigma_j\), it is the same coefficient-defined derivative on
the smaller strip, by the compatibility following
\Cref{ana:wiener-estimate}.  Therefore
\[
\begin{split}
|C(u_{j+1})|_{\sigma_{j+1}}
&\le
C\|\eta_j\|_{r_0,\sigma_{j+1}}
 |\eta_j|_{\sigma_{j+1}}\\
&\le C P^{r_0}2^{r_0j}\eps_j^2.
\end{split}
\]
Equivalently,
\begin{equation}\label{ana:newton-recurrence}
\eps_{j+1}\le
C P^{r_0}2^{r_0j}\eps_j^2.
\end{equation}
This calculation can first be performed coefficientwise on the real
axis, where every variation is smooth, and then estimated in the
smaller-strip norm.  It does not assume differentiability on a
Wiener space with insufficient boundary derivatives.

Choose \(P\) sufficiently large that
\(\eps_0\le P^{-2r_0}\), as permitted by
\Cref{ana:initial-residual}.  Increase it further so that the fixed
constant in \Cref{ana:newton-recurrence} satisfies
\(CP^{-r_0}2^{r_0+1}<1\).  Induction gives
\begin{equation}\label{ana:newton-decay}
\eps_j\le P^{-2r_0}2^{-(r_0+1)j}.
\end{equation}
Indeed the ratio of the right hand side of
\Cref{ana:newton-recurrence}, with this bound substituted, to the
proposed bound at \(j+1\), is at most
\(CP^{-r_0}2^{r_0+1}2^{-j}\).
Moreover, \Cref{ana:correction-smoothing,ana:newton-decay} give
\[
\sum_{j\ge0}\|\eta_j\|_{r_0,\sigma_{j+1}}
\le CP^{-r_0}\sum_{j\ge0}2^{-j}.
\]
For sufficiently large \(P\) this is smaller than \(M/4\).
By \Cref{ana:truncation-bound}, \(u_0\) has norm less than \(M/4\).
Every iterate and every correcting segment therefore has the required
norm with slack.  This proves the inductive bounded-input condition
rather than assuming it.
Formally, the induction is simultaneous: before applying
\Cref{ana:newton-recurrence} at step \(j\), the already known bound
on \(\eps_j\) and \Cref{ana:correction-smoothing} put that entire
correcting segment within the same cumulative norm bound; the
recurrence then proves the next error bound.

The corrections converge in \(C^{r_0}(\mathcal A_{\sigma_\infty})\)
to a real-symmetric holomorphic function \(u_*\).
Its low Fourier modes agree with those of \(u\).  For each fixed
constraint, continuity of the real mesh construction and
\(\eps_j\to0\) give \(C_n(u_*)=0\).
The real segment from \(u\) to \(u_*\) remains in the original bounded
input set.  Put \(v=u_*-u\), which has only high Fourier modes.
At width zero,
\[
0=C(u_*)-C(u)
=v+\int_0^1R(u+\lambda v)v\,\dd\lambda .
\]
The operator norm of the integral is at most \(1/2\).  Hence
\(|v|_0\le|v|_0/2\), so \(v=0\).
Thus the original \(u\) is real analytic.

By \Cref{ana:normalization}, \(\phi'=e^u>0\) on the real axis.
The primitive \(\phi\) and its inverse are real analytic.
The same equation then gives analytic \(r\) in the normal-angle
coordinate.  Finally the smooth support function solves
\(h''+h=r\), an analytic ordinary differential equation, and is
therefore analytic.  Its boundary parametrization
\(h n_\theta+h'e_\theta\) is analytic and regular.
\end{proof}

\begin{corollary}\label{ana:analytic-action-interface}
For the actual table, the normalized action in
\Cref{ana:normalized-action} is holomorphic in a neighborhood of the
real diagonal, diagonally \(2\pi\)-periodic, and antisymmetric.
It satisfies
\[
D(X,X+T)=-T^3/24+O(T^5),\qquad
D_1(X,X+T)=T^2/8+O(T^4),
\]
locally uniformly with all fixed derivatives.  The unique rational
graphs and their identities in \Cref{lem:rational-periodicity}
are available for every sufficiently small positive reduced \(l/q\).
\end{corollary}

\begin{proof}
Holomorphic extension follows from the analytic \(u\), the midpoint
implicit equation, and the integral formula.
The identities and expansions are
\Cref{ana:action-symmetries,ana:action-expansion};
the final assertion is \Cref{lem:rational-periodicity}.
\end{proof}
